% M08_c10_campaign.tex -- the shared-model ATLAS campaign C10 (decision D-13): design, registered hypotheses H-A..H-F,
% H-J, the S0/S1 results that were seen before register revision 2, and the sealed state.
% Sources: analysis_2026-10-10_campaign/CAMPAIGN_DESIGN.md (sections 1-7), DFT_TO_TCAD.md, S2_PREP.md;
% HYPOTHESIS_REGISTER.md section 0 (S0/S1 results "seen before revision 2") and H-A..H-F, H-J, A13, A41; decision log
% D-13, D-17, D-44, D-58. SIGHT RULE: no outcome of S2 or later appears here (sealed runs 0075-0077, campaign
% C10_DIAGb; commit messages that mention S2 scores are treated as sealed-adjacent and not used).
\chapter{The shared-model device-simulation campaign (C10)}
\label{ch:mono-c10}

The V1 and V2 models of \cref{ch:calibration,ch:v2} fitted a threshold-setting charge and a mobility to each film, and
every out-of-sample threshold test at 6.3~nm failed (\claim{C014}). On 10 October 2026 the user approved a campaign of
up to 200 further ATLAS launches with a different aim (\dlog{D-13}): one parameter set shared by the 2.0, 6.3 and
13.2~nm films, first-principles quantities entering as computed values or as laws fitted to computed points, explicit
tests of the off-state floor, and every fitted versus computed parameter labelled. This chapter describes the design,
the hypotheses registered for it before their outcomes, the two stages whose results were seen before the register was
frozen, and the present state, in which every later outcome is sealed.

% =====================================================================================================================
\section{Design principles}
\label{sec:mono-c10-design}

\paragraph{One shared model.} Five fitted parameters are shared by all films: the alignment offset $\Delta V_0$ (carried
by a front fixed charge $\Qf$), the bulk mobility $\mu_\infty$, a scattering length $\ell$, the integrated tail density
$N_s$ (with $N_t = N_s/t$) and the tail width $W_T$; a sixth only if a percolation-mobility branch is adopted. No
film-specific charge or mobility is allowed. First-principles inputs enter through the laws of
\cref{tab:mono-mass-laws}: $\dEc(t)$ in the affinity, $\Delta\Eg(t)$ in the gap, $\Delta m^*(t)$ in $\Nc$ and in the
mobility law, and a quasi-two-dimensional $\Nc$ at 2.0~nm ($1.373\times10^{19}$\cmcu{}, \cref{sec:mono-conf-dos}). The
mobility of each film follows a Matthiessen form,
\begin{equation}
\mu_{\mathrm{band}}(t) = \mu_\infty\,\frac{0.208\,m_0}{m^*(t)}\,\frac{1}{1+\ell/t}\,\Big(1-\frac{0.287~\mathrm{nm}}{t}\Big)^2 ,
\label{eq:mono-c10-mu}
\end{equation}
in which the mass ratio (literature bulk mass with the DFT increment) explains only a factor 0.83 of the drop from 13.2
to 2.0~nm and the fitted $\ell$ must carry the rest; the last factor is the existing roughness term.

\paragraph{Exact offline profiling.} Two of the shared parameters can be fitted without launches, because the
drift--diffusion equations transform exactly. A uniform front sheet charge changes only the oxide field, so the transfer
curve shifts rigidly, $\Id(\Vg;\Qf) = \Id(\Vg + q(\Qf-\Qf^{(0)})/\Cox;\Qf^{(0)})$; and with an electron-only model in
which every mobility is multiplied by one constant, $\Id \propto$ MUN. The generation--recombination term breaks the
second identity, but it is at most $10^{-28}$~A/\textmu m against currents of at least $10^{-15}$~A/\textmu m
(calculated bound). $\Delta V_0$, $\mu_\infty$ and $\ell$ are therefore profiled offline for every launched film triple,
leaving three degrees of freedom as the honest test of sharing; every optimum is confirmed natively.

\paragraph{Objective and acceptance, fixed before any launch.} The objective weights the films equally,
$J = \tfrac13\sum_t\tfrac12(\mathrm{RMSE}^2_{\mathrm{sub},t} + \mathrm{RMSE}^2_{\mathrm{on},t})$, with log-current residuals
in the subthreshold window (measured $\Id$ between ten times the floor and $10^{-9}$~A/\textmu m) and over
$\Vg = 1$--3~V. Acceptance per film (assumed engineering limits): active RMSE $\le0.15$~dec, $|\Delta\Vthcc|\le0.10$~V,
$\Id(3~\mathrm{V})$ within $\pm30$\,\%, SS (registered window) within $\pm20$\,\%, $\muFE$ at 1.5, 2 and 3~V within
$\pm20$\,\%. A film that fails any limit is reported as failed, and success on all three is still called calibration.

\paragraph{Expected outcome, stated in the design.} The design states that the shared model will probably not fit the
6.3~nm film: its on-current is \emph{below} that of the 2.0~nm film, which neither a monotonic mobility law nor a
monotonic first-principles shift can produce. No per-film rescue parameter is allowed; instead the residual is split
into a rigid shift and a shape part.

\begin{table}[htbp]\centering\small
\caption{Model set of the campaign (abridged; full list and reasons in \file{CAMPAIGN_DESIGN.md} section 3).}
\label{tab:mono-c10-models}
\begin{tabularx}{\linewidth}{@{}>{\raggedright\arraybackslash}p{5.3cm}>{\raggedright\arraybackslash}p{2.1cm}>{\raggedright\arraybackslash}X@{}}
\toprule
Model & Status & Reason \\\midrule
Fermi--Dirac statistics & primary & degenerate accumulation ($\approx5.6\times10^{19}$\cmcu{} at 2~nm, 2~V overdrive) \\
DFT confinement in affinity, gap, $\Nc$; quasi-2-D $\Nc$ at 2~nm & primary & computed points, fitted laws (fit3) \\
Exponential acceptor tail, $N_t = N_s/t$ & primary, fitted & swing collapse in data and TCAD \\
Deep Gaussian; interface acceptor sheet & on (assumed) & continuity with V1/V2 \\
Uniform donors $2.5\times10^{17}$\cmcu{} & on (V1 fitted amplitude, fixed) & compensation not computed; thickness term retired \\
Front fixed charge & primary, fitted & carries $\Delta V_0$ (degenerate with work function and affinity) \\
Shared mobility law, \cref{eq:mono-c10-mu} & primary, fitted & Matthiessen form \\
IGZO.TOKYO percolation mobility & branch (6 runs) & targets the rising $\muFE$ and $\gamma$ \\
Schottky Pd contacts, shared barrier & branch (7 runs) & on-state shape \\
Bohm quantum potential & test (2 runs) & never together with the DFT shift \\
SRH, trap-assisted tunnelling, two carriers & test & quantitative null tests of generation \\
Gate tunnelling (self-consistent), Fowler--Nordheim, oxide TAT & test & gate leakage as a floor candidate \\
Ungated strip (separate deck, in-solver) & test & the only candidate that gives a $\Vg$-flat floor of the right size \\
Field-dependent, silicon inversion-layer and doping mobility models; band-to-band tunnelling; lumped contact resistance & off & physically inapplicable or
overwrite the material mobility \\\bottomrule
\end{tabularx}
\end{table}

\section{Staged run matrix and stop rules}
\label{sec:mono-c10-matrix}

\begin{table}[htbp]\centering\small
\caption{Planned stages (181 launches + 19 reserve = 200; the regenerated matrix has 184 entries and 16 reserve, a
recorded inconsistency, \cref{app:mono-corrections}).}
\label{tab:mono-c10-stages}
\begin{tabularx}{\linewidth}{@{}lr>{\raggedright\arraybackslash}X@{}}
\toprule
Stage & Launches & Content \\\midrule
S0 smoke & 2 & 2.0~nm primary baseline; the same with worst-case gate tunnelling \\
S1 baselines and numerics & 7 & 6.3 and 13.2~nm baselines; mesh and DOS refinement; 3-D $\Nc$ at 2.0~nm; two carriers; native profiling check \\
S2 off-floor mechanisms & 17 & generation bounds; gate leakage (nominal, worst case, FN, oxide TAT, mesh check); ungated-strip decks and controls \\
S3 model-form branches & 27 & confinement off; tail placement; IGZO; Schottky; Bohm; Boltzmann; native check \\
S4 calibration & 69 & 45 space-filling (Latin hypercube, seeds fixed), 18 adaptive, 3 native confirmation, 3 runner-up \\
S5 leave-one-film-out & 30 & three folds; the held-out film is never tuned \\
S6 sensitivity and frozen predictions & 27 & confinement low/HSE-high, bulk mass, $N_D$, $\varepsilon$(\AlO), traps off, 6.3~nm quasi-2-D, frozen strip decks \\
S7 31.8~nm stress & 2 & frozen model; documented, never fitted \\\bottomrule
\end{tabularx}
\end{table}

The stop rules are fixed in the design: numerics refinements are adopted if they change the active RMSE by more than
0.01~dec; a branch is adopted only if it lowers $\sqrt J$ by at least 0.03~dec without worsening any film by more than
0.03~dec; if after S4a the best $\sqrt J$ exceeds 0.30~dec with systematic residuals at 6.3~nm, the shared model is
reported as failed, with no per-film rescue; and if S5 confirms failure on two or more folds, the sensitivity stage is
cut to the first-principles input profiles. Every launch counts, failures included; the cap of 200 is enforced in code;
runs are strictly sequential.

% =====================================================================================================================
\section{Registered hypotheses}
\label{sec:mono-c10-hyp}

The hypotheses were frozen in register revision~2 on 2026-10-10T18:33:13Z (\dlog{D-17}), after the S0 and S1 results
were seen and before any S2 or later launch. \Cref{tab:mono-c10-hyp} summarises the registered predictions; the
register gives the scoring rules, the censoring rule (a quantity below the demonstrated numerical resolution counts as
consistent with an upper bound only if the resolution is below the bound; otherwise unresolved) and the forecasts,
which are recorded but not scored.

\begin{table}[htbp]\centering\footnotesize
\caption{Registered device-model hypotheses of campaign C10 (abridged). Outcomes: \sealedph{S2 and later}.}
\label{tab:mono-c10-hyp}
\begin{tabularx}{\linewidth}{@{}>{\raggedright\arraybackslash}p{0.8cm}>{\raggedright\arraybackslash}p{3.6cm}>{\raggedright\arraybackslash}X@{}}
\toprule
ID & Hypothesis & Registered prediction \\\midrule
H-A & the 2.0~nm floor is extrinsic or instrument-limited & converged intrinsic mechanisms stay $\ge1$~dec below
$4.6\times10^{-15}$~A/\textmu m, with resolution shown below that bound; decisive part expected unresolved (no instrument records) \\
H-B & an ungated strip sets the floors & $I(0.7~\mathrm{V})/I(0.1~\mathrm{V}) = 7.0\pm0.5$; with $(W/L)_{\mathrm{par}}$
fitted on 13.2~nm, the 6.3 and 2.0~nm floors within 0.5~dec \\
H-C & gate leakage and generation do not set the floor & nominal leakage $\le10^{-2}\times$ floor (if resolved);
extreme generation $\le10^{-20}$~A/\textmu m; worst-case leakage fails flatness or thickness scaling \\
H-D & band mobility needs $1/t$ scattering & scattering branch gives a 2.0/13.2~nm band-mobility ratio of 0.20--0.35
with $\ell = 4$--12~nm (mass-only ratio unresolved: amorphous dispersion workflow descoped, \amend{A22}) \\
H-E & confinement and the threshold & all inputs off: $\Delta\Vthcc = -0.13\pm0.03$~V (2.0~nm; revision-1 value $-0.23\pm0.04$
kept and scored), $-0.062\pm0.015$ (6.3), $-0.023\pm0.008$~V (13.2); affinity-only runs rigid within $\pm5$~mV \\
H-F & the 6.3~nm shape residual is unexplained & two-film reference needs a rigid shift of $+0.28\pm0.08$~V and a
current scale of 0.15--0.35; baseline shape RMSE $>0.15$~dec \\
H-J & one shared model transfers across 2.0, 6.3, 13.2~nm & acceptance of \cref{sec:mono-c10-design} for every film in
every fold; any failure falsifies transfer \\\bottomrule
\end{tabularx}
\end{table}

\paragraph{Forecasts.} Recorded before the freeze: Claude forecast that transfer (H-J) fails, with a best space-filling
$\sqrt J$ of 0.25--0.65~dec and, with 6.3~nm left out, a threshold miss of $-0.33\pm0.10$~V; Codex first forecast
success and withdrew that forecast as overconfident on 11 October. For H-D, the mass alone is forecast to explain about
14\,\% of the logarithmic mobility drop, with $\ell = 8\pm3$~nm. These are forecasts, not tests.

% =====================================================================================================================
\section{Results seen before the register freeze (S0, S1)}
\label{sec:mono-c10-s1}

Launches run\_0061 to run\_0074 (14 of the 200) were made and seen before revision 2 was frozen; the register lists their
outcomes as known, and they are therefore reported here (\cref{tab:mono-c10-s1}). They are model-internal results, at the
shared start parameters, not calibrations.

\begin{table}[htbp]\centering\small
\caption{Stages S0 and S1 (register section 0; runs 0061--0074). RMSE: active log-current residual against the measured
curve at the shared start parameters.}
\label{tab:mono-c10-s1}
\begin{tabularx}{\linewidth}{@{}>{\raggedright\arraybackslash}p{1.5cm}>{\raggedright\arraybackslash}p{4.4cm}>{\raggedright\arraybackslash}X@{}}
\toprule
Run & What it tested & Result \\\midrule
S0-01 & 2.0~nm primary baseline & deck parameters reproduced in the deck header (basis for later runs) \\
S0-02 & 2.0~nm with worst-case gate tunnelling & needed six launches against a cap of three repairs (disclosed protocol deviation) \\
S1-01 & 6.3~nm baseline & RMSE 1.18~dec; $\Id(3~\mathrm{V})$ $\times4.87$ the measured value \\
S1-02 & 13.2~nm baseline & RMSE 0.056~dec \\
S1-03 & 13.2~nm, mesh $\times0.7$ & passes the numerics rule ($\Delta$RMSE $-0.0005$~dec) \\
S1-04 & 2.0~nm, DOS quadrature doubled & passes ($-0.00002$~dec) \\
S1-05 & 2.0~nm, 3-D instead of quasi-2-D $\Nc$ & $\Delta\Vthcc = +142$~mV, outside the registered $+40$ to $+70$~mV:
\textbf{failed}; $\Delta$SS $+49$~mV/dec; RMSE $0.367\to0.053$~dec \\
S1-06 & 13.2~nm, two carriers & timed out: unresolved \\
S1-07 & native check of the offline profiling & passed: largest deviation $2.05\times10^{-4}$~dec over both windows
(validates the tested $\Qf$/MUN transformation only) \\\bottomrule
\end{tabularx}
\end{table}

\paragraph{What S1 already shows.} The 6.3~nm baseline at the shared start parameters is off by more than a decade
and overestimates the on-current almost five-fold, as the design anticipated; the 13.2~nm baseline is close. The failed
S1-05 expectation was audited (\amend{A13}) and no implementation error was found: of the $+142$~mV, $+37.5$~mV come from
the statistics change and $+89.5$~mV from additional charge in the 40~meV acceptor tail, because at $10^{-9}$~A/\textmu m
the effective slope factor is about 3.5, not the 1.47 the registration assumed. The audit also showed that, at
2.0~nm, the quasi-2-D $\Nc$ approximation differs from true two-dimensional statistics at matched sheet density by only
0.86~mV at threshold but overstates the on-current by about 7--8\,\% (inferred), a disclosed model limitation. Because
$\Nc$ cannot be separated from $\Qf$, $N_{TA}$ or the Gaussian depth using subthreshold transfer data, every fitted trap
density and fixed charge must state its $\Nc$ convention. The H-E prediction at 2.0~nm was re-registered after S1 and
before S3 ($-0.13\pm0.03$~V), with the original revision-1 value kept and scored as well.

% =====================================================================================================================
\section{Present state: sealed, budget and conditions}
\label{sec:mono-c10-state}

After the freeze, the campaign was prepared for S2 with a frozen scorer (\texttt{score\_register.py},
sha256 \texttt{c0f28bb9}\ldots), a pre-launch gate and runner rules for state-only reporting of sealed runs
(\amend{A41}, \dlog{D-44}). Three later launches (run\_0075 to run\_0077) and one diagnostic campaign are sealed under
that rule; they are absent from this repository, and the monograph contains no outcome of them. By 14 October, 17 of the
200 launches had been used. The user accepted a plan for the remaining campaign (\dlog{D-58}): a two-carrier repair of
S1-06 first, the remaining diagnostic set (16 launches) and two further gate-leakage runs, then the tunnelling-mesh and
Fowler--Nordheim branches (6 launches), each conditional on the independent reviewer's approval of its amendment and on
the pre-launch gate; sensitivity runs are cut first if launches run short.

\paragraph{What the campaign can establish.} Within the model: whether one parameter set transfers across the films
(H-J), whether the off-state floors can be produced by gate leakage, generation or an ungated strip (H-A to H-C), and how
much of the threshold shift the first-principles inputs carry (H-E). It cannot establish the device physics by itself:
four single-device curves, one drain bias and no gate current leave several mechanisms indistinguishable, and the
confinement inputs are not identifiable from transfer curves (\claim{C018}). The decisive measurements are listed in
\cref{ch:mono-conclusions}.

\begin{keyresult}
Campaign C10 replaces per-film fitting by one shared parameter set with first-principles laws, exact offline profiling
of the shift and scale, acceptance criteria and stop rules fixed before any launch, and a stated expectation that the
6.3~nm film will not be fitted. The two stages seen before the register freeze show the 6.3~nm baseline off by more than
a decade and a failed registered expectation for the density-of-states convention, explained by tail-state charging.
All later outcomes, including every test of the off-state floor and the transfer test, are sealed at the evidence
state of this monograph.
\end{keyresult}
